\section{評価と考察}

4.2節で示した結果について評価及び考察を行う．

\subsection{覚醒時と覚醒度低下時のRRI変動の違いについて}

4.2.2節で述べたように図\ref{fig15}に示すRRIの値は図\ref{fig11}に比べ，安定している．逆に図\ref{fig11}に示すRRIの値は図\ref{fig15}に比べ，変動している．つまり，覚醒度低下時と覚醒時では，RRI変動に違いが見られる．この理由について考える．
呼吸周期が長くなると，RRIの変化が呼吸性変動に対して追随しやすくなると言われている\cite{[11]}．覚醒度が低下すると，副交感神経が働くことにより，呼吸周期は長くなり，呼吸パターンが安定する．よって，安定した呼吸パターンにRRIが追随するため，覚醒度低下時にRRIの値が安定すると考えられる．

一方，呼吸周期が短いと，RRIの変化が呼吸性変動に対して追随しにくいと言われている\cite{[11]}．覚醒時には，交感神経が活発化し，ストレスや活動量などにより呼吸周期が短くなり，呼吸パターンが変動する．よって，RRIが安定したパターンに追随することがないため，覚醒時にRRIの値が変動すると考えられる．

\subsection{本実験で用いたLPの信頼性について}

2.2.2節で述べたように，安静時，LPのプロットの重心が右上に推移し，緊張時は左下方向に推移しながら各点の散らばりが大きくなる．さらに緊張が高まると，プロットの重心は左下に推移しながら各点の散らばりが小さくなることが知られている\cite{[9]}．

4.2.1節で図\ref{fig10}を用いて示したように，覚醒時のLPの各点には， 200msから1000msまでのばらつきが見られた．また，4.2.2節で図\ref{fig14}を用いて示したように，覚醒度低下時のLPの各点は，$RRI_n-RRI_{n+1}$平面上において右上に720msから900msの値の範囲に集まってプロットされた．これらは，2.2.2節で述べた理論通りの結果となっている．よって，本実験で実装したLPは，被験者の覚醒度を幾何学的に正しく示していると言える．

\subsection{F-stressの技術的課題について}

2.2.3節では，F-stressがゼロに近くなるほど，ストレスが高まっている\cite{[9]}と述べた．しかし，覚醒度低下時にもかかわらず，4.2.2節で述べた図\ref{fig16}に示すF-stressの値は，0秒から50秒にかけて0に近い値を取っている．この理由を考察する．

2.2.3節で示したF-stressを求める計算式に注目する．(3)式で示しているようにF-stressは$g_p$と$\delta$の積であった．ところで，$\delta$は各点から重心までの平均距離であり，重心と各点のばらつきを表す指標であった．覚醒度低下時は，5.1節で述べたようにRRI変動は小さくなり，図\ref{fig14}のようにLP上の各点のばらつきは小さくなる．よって，覚醒度低下時の$\delta$の値が0に近づくため，F-stressも図\ref{fig16}のように0に近い値を取るのだと考えられる．

確かに，2.2.2節で述べたように緊張度がかなり高いと$RRI_n-RRI_{n+1}$平面上において左下にプロットが集中することが知られている\cite{[9]}．しかし，覚醒度低下時にも散らばりが小さくなるため，どちらの場合においてもF-stressは0に近い値を取ることになる．よって，評価指標F-stressは，ストレスが高い時と覚醒度低下時の区別ができないという技術的課題が，本実験の結果より明らかになった．


\subsection{本装置の居眠り検知・防止の評価}

起床時の結果を評価する．4.2.3節で述べたように，図\ref{fig18}に示す起床時のLPは，図\ref{fig14}に示す覚醒度低下時のLPに比べ，各点がばらつき始めている．具体的には， 720msから900msの間にあったプロットが200msから900msの間に散らばり始めている．2.2.2節で述べた理論から，本装置によって，安静時から被験者の緊張が徐々に高まっていると考えられる．また，4.2.3節で述べたように，被験者はモータの振動によって30秒から31秒の間に目覚めている．実際に起床したという事実と5.2節で述べた信頼性のある幾何学的証拠から，本装置は居眠りの検知・防止に有効性があると言える．

\subsection{F-stressに対する覚醒度解析における有効性}

4.2.3で述べたように，図\ref{fig20}に示すF-stressは，0秒から30.5秒においても0に近い値を取っていた．しかし，5.4節で述べたようにF-stressは覚醒度低下時とストレスが高い時を区別できない．一方，図\ref{fig21}に示す光センサBのグラフは0秒から25秒にかけて周期的な波形をしていた．4.2.2節で示した図\ref{fig17}の光センサBのグラフも同様に周期的な波形をしていた．どちらも被験者がうとうとしていたことから，この周期的な波形は，頭のうとうととした動きを捉えていると考えられる．本システムは，このうとうとしている様子をグラフ上に捉えられるか否かにおいて，既存の覚醒度解析よりも有効性があると言える．

\subsection{誤使用時，暗所にいる際の動作の評価}

4.2.4節の結果から，誤使用時の動作が2.1.1節で示した通りに動いていることがわかる．具体的には，光センサBの値＞光センサAの値，光センサAの値＞光センサBの値と切り替えることによって，正しくStateが$0\rightarrow 2$，$2\rightarrow 0$と切り替わっていることがわかる．また，モータの動作も同じように停止$\rightarrow$回転，回転$\rightarrow$停止と切替わっていることがわかる．よって，誤使用時のシステムは本装置に正しく実装できたと言える．

また同様に，4.2.5節の結果から，暗所にいる際の動作が2.1.1節で示した通りに動いていることが分かる．具体的には，部屋の照明を点灯$\rightarrow$消灯，消灯$\rightarrow$点灯と切り替えることによって，正しくModeが$0\rightarrow 1$，$1\rightarrow 0$に切替わっていることがわかる．また，暗所において，モータが回らないことも確認できた．よって，暗所のシステムが本装置に正しく実装できたと言える．

